\section{Results}
\label{sec:results}

\subsection{A 2.7M-parameter model against foundation models}
\label{sec:main}

Table~\ref{tab:main} compares CroFreMo against the strongest reproduced
baseline on each corpus; the full matrix of fifteen baselines is in the
appendix. The axis-free
model wins on the two seizure-detection corpora with adequate positives
(TUSZ, CHB-MIT) by margins far outside the baselines' seed variance, wins
narrowly on IIIC and TUEP, ties the strongest foundation model on TUEV and
the strongest supervised model on TUAB, and loses on the two sleep-staging
corpora, on the two cohort-diagnosis corpora, on TUAR and on Siena. On
CHB-MIT the strongest baseline is not the largest model: CBraMod trained
from scratch reaches $0.317 \pm 0.205$ and its released pretrained weights
$0.233 \pm 0.098$, against $0.716$ for a model thirteen times smaller.


\begin{table}[t]
\centering
\small
\caption{CroFreMo (2.7M, trained from scratch) against the strongest
reproduced baseline per corpus. Baselines are three-seed means $\pm$ std
unless marked $(1)$; CroFreMo is a single seed throughout. ``pre'' marks a baseline evaluated with its released
pretrained weights. TUAB is scored by balanced accuracy for every model.}
\label{tab:main}
\resizebox{\textwidth}{!}{%
\begin{tabular}{lllcc}
\toprule
Corpus & Metric & Strongest baseline & & CroFreMo \\
\midrule
TUSZ & AUC-PR & FFCL & $0.545 \pm 0.030$ & $\mathbf{0.698\,(1)}$ \\
CHB-MIT & AUC-PR & TFM-Tokenizer pre & $0.627 \pm 0.025$ & $\mathbf{0.716\,(1)}$ \\
IIIC & $\kappa$ & REVE pre & $0.436 \pm 0.003$ & $0.447\,(1)$ \\
TUEP & AUROC & EEGPT scratch & $0.786 \pm 0.021$ & $0.799\,(1)$ \\
TUEV & $\kappa$ & REVE pre & $0.685 \pm 0.039$ & $0.680\,(1)$ \\
TUAB & bal.\ acc. & ST-Transformer & $\mathbf{0.820 \pm 0.004}$ & $0.820\,(1)$ \\
Sleep-EDF & $\kappa$ & ContraWR & $\mathbf{0.692 \pm 0.015}$ & $0.639\,(1)$ \\
CAUEEG & bal.\ acc. & BIOT scratch & $\mathbf{0.561 \pm 0.012}$ & $0.511\,(1)$ \\
ISRUC & $\kappa$ & CBraMod pre & $\mathbf{0.754 \pm 0.008}$ & $0.698\,(1)$ \\
ADFD & bal.\ acc. & BIOT scratch & $\mathbf{0.525 \pm 0.021}$ & $0.461\,(1)$ \\
TUAR & $\kappa$ & CBraMod pre & $\mathbf{0.715^{(1)}}$ & $0.600\,(1)$ \\
Siena & AUC-PR & REVE pre & $\mathbf{0.518 \pm 0.117}$ & $0.201\,(1)$ \\
\bottomrule
\end{tabular}}
\end{table}

\subsection{What the seizure margin is made of}
\label{sec:attribution}

Table~\ref{tab:attribution} is the experiment the encoder was built for.
The axis-free model is trained with its tokenizer producing either the
duplex grid of Section~\ref{sec:tokenizer} or waveform tokens only, with
every other setting identical, at the adopted width and at one width
larger. Compared with the waveform-only tokenizer, the duplex tokenizer improves
CHB-MIT by $0.136$ AUC-PR, TUEV by $0.107$ $\kappa$, TUSZ by $0.063$ AUC-PR
and IIIC by $0.027$ $\kappa$ at $d{=}192$. At $d{=}256$ the four effects keep their sign; three
are of similar size and the CHB-MIT effect drops to $0.025$, so the size
of the effect is sensitive to capacity and optimization state, not only
to the corpus. For scale, the seed standard deviations of the same
recipe on the same corpora, measured on an earlier design with a
dedicated frequency axis (Appendix~\ref{app:triaxial}), are $0.055$, $0.036$, $0.035$ and $0.009$; they
are a reference for the recipe, not a substitute for this model's own
variance. The two margins by which the model leads the
foundation-model field are therefore not properties of the band
decomposition or of the folded attention alone: without the interaction
rows the same encoder falls to $0.635$ on TUSZ and $0.580$ on CHB-MIT,
and on CHB-MIT below TFM-Tokenizer.

Table~\ref{tab:perclass} breaks the TUEV effect down by class. Every
class improves, and the gain is largest on the two periodic-discharge
classes (GPED F1 $0.59 \to 0.83$, PLED $0.50 \to 0.58$), whose defining
morphology is fast activity locked to a slow rhythm; the rare
spike-and-sharp-wave class moves from near-zero recall to $0.16$, and
background, artifact and eye movement change little. The confusion that
remains is between GPED and PLED and between artifact and background.

This comparison changes more than one thing at once --- the waveform-only
grid has half the rows, no explicit envelope feature and no phase
representation --- so it attributes the gain to the interaction rows as a
whole, not yet to the cross-frequency alignment inside them. Controls
that hold the row count fixed and vary what the extra rows carry ---
waveform rows, amplitude-only rows, analytic features without cross-band
alignment --- are the experiment that would separate these, and they are
not in this version.

\begin{table}[t]
\centering
\small
\caption{Duplex vs.\ waveform-only tokenizer in the axis-free encoder,
everything else identical. Single seeds.}
\label{tab:attribution}
\begin{tabular}{llcccccc}
\toprule
 & & \multicolumn{3}{c}{$d{=}192$ (final)} & \multicolumn{3}{c}{$d{=}256$} \\
\cmidrule(lr){3-5}\cmidrule(lr){6-8}
Corpus & Metric & duplex & waveform-only & $\Delta$ & duplex & waveform-only & $\Delta$ \\
\midrule
CHB-MIT & AUC-PR & $\mathbf{0.716}$ & $0.580$ & $\mathbf{+0.136}$ & $0.608$ & $0.583$ & $+0.025$ \\
TUEV & $\kappa$ & $\mathbf{0.680}$ & $0.573$ & $\mathbf{+0.107}$ & $0.679$ & $0.601$ & $+0.078$ \\
TUSZ & AUC-PR & $\mathbf{0.698}$ & $0.635$ & $\mathbf{+0.063}$ & $0.631$ & $0.541$ & $+0.090$ \\
IIIC & $\kappa$ & $\mathbf{0.447}$ & $0.420$ & $+0.027$ & $0.465$ & $0.418$ & $+0.047$ \\
\bottomrule
\end{tabular}
\end{table}

\begin{table}[t]
\centering
\small
\caption{TUEV per class, CroFreMo with the waveform-only vs.\ the duplex
tokenizer (single seed each, 29,421 test windows): test-set count and
precision / recall / F1. Macro-F1 $0.48 \to 0.58$.}
\label{tab:perclass}
\begin{tabular}{lrccc}
\toprule
Class & $n$ & waveform-only P / R / F1 & duplex P / R / F1 & $\Delta$F1 \\
\midrule
SPSW & 567 & 0.87 / 0.04 / 0.07 & 0.29 / 0.16 / 0.20 & $+0.13$ \\
GPED & 4,677 & 0.61 / 0.56 / 0.59 & 0.82 / 0.83 / 0.83 & $+0.24$ \\
PLED & 1,998 & 0.48 / 0.53 / 0.50 & 0.63 / 0.54 / 0.58 & $+0.08$ \\
EYEM & 329 & 0.94 / 0.23 / 0.37 & 0.44 / 0.42 / 0.43 & $+0.06$ \\
ARTF & 2,204 & 0.35 / 0.62 / 0.44 & 0.42 / 0.65 / 0.51 & $+0.07$ \\
BCKG & 19,646 & 0.93 / 0.89 / 0.91 & 0.94 / 0.91 / 0.92 & $+0.01$ \\
\bottomrule
\end{tabular}
\end{table}


\subsection{The same token in other encoders}
\label{sec:content-structure}

Table~\ref{tab:content-structure} collects the tokenizer's effect across
four encoders. The interventions are not the same experiment, and the
table says which is which: for CroFreMo it is the duplex tokenizer against
a waveform-only one; for CBraMod it is interaction rows against matched
waveform rows added to the native tokenizer; for LaBraM and REVE it is the
whole CroFreMo tokenizer, band decomposition included, against the host's
own. CBraMod's encoder, which keeps its own rFFT magnitude branch, gains from
the interaction rows relative to matched waveform rows on TUEV ($0.045$),
TUEP ($0.057$), IIIC ($0.015$) and CHB-MIT ($0.124$), and loses $0.092$ on
TUSZ. LaBraM and REVE, trained from scratch with no spectral pathway,
gain $0.150$ and $0.241$ on TUEV when the tokenizer replaces their patch
embedding. CroFreMo gains on all four corpora measured.

\begin{table}[t]
\centering
\small
\caption{Tokenizer benefit across four encoders. The interventions differ
and the numbers are not one quantity: duplex vs.\ waveform-only tokenizer
(CroFreMo); interaction rows vs.\ matched waveform rows added to the
native tokenizer (CBraMod); the whole CroFreMo tokenizer vs.\ the native
tokenizer (LaBraM, REVE). Single seeds.}
\label{tab:content-structure}
\resizebox{\textwidth}{!}{%
\begin{tabular}{lllccccc}
\toprule
Encoder & Intervention & Control & TUEV & TUSZ & CHB-MIT & IIIC & TUEP \\
\midrule
CroFreMo & duplex tokenizer & waveform-only tokenizer & $+0.107$ & $+0.063$ & $+0.136$ & $+0.027$ & --- \\
CBraMod (scratch) & + interaction rows & + waveform rows & $+0.045$ & $-0.092$ & $+0.124$ & $+0.015$ & $+0.057$ \\
LaBraM (scratch) & CroFreMo tokenizer & native tokenizer & $+0.150$ & --- & $+0.017$ & $+0.014$ & --- \\
REVE (scratch) & CroFreMo tokenizer & native tokenizer & $+0.241$ & --- & --- & $+0.078$ & --- \\
\bottomrule
\end{tabular}}
\end{table}

Read within each intervention, the pattern is this. With its own
tokenizer replaced, LaBraM and REVE gain most on TUEV; that gain includes
the band decomposition and cannot yet be split between it and the
interaction rows, which would need the CroFreMo waveform-only tokenizer
in the same hosts. With the native tokenizer kept and rows added, CBraMod
gains on the three event-classification corpora and on CHB-MIT, where it
trains poorly from scratch and the interaction rows give it a foothold,
and loses on TUSZ. Whether an encoder with a spectral branch computes for
itself what the interaction rows supply is a question for probing or
intervention experiments; the table records that the benefit of explicit
cross-frequency content varies with the encoder and the corpus, not why.

\subsection{Design controls}
\label{sec:design}

Table~\ref{tab:design} reports the variants that were tried on the way
to the final model, all single seeds. Folding the band rows
into spatial attention rather than leaving them as unrelated channels is
worth $0.09$ on both seizure corpora and costs $0.02$--$0.04$ on the
smaller corpora, where the longer attention span overfits. Raising the
width from 192 to 256 recovers those small-corpus cells and loses
$0.07$ and $0.11$ on TUSZ and CHB-MIT. An explicit coupling-strength
feature --- the magnitudes $|Z_{ij}|$ projected onto the fused rows ---
raises TUEV to $0.721$ and Siena to $0.398$ and collapses CHB-MIT to
$0.513$. No variant dominates; the adopted one is the one that keeps the
seizure margin, and every alternative is reported. Two controls run on an earlier design with a dedicated frequency axis,
reported in Appendix~\ref{app:triaxial}, locate where a small model's
seizure-detection margin comes from: collapsing the learned filterbank to
one broadband band costs $0.11$ AUC-PR on TUSZ and $0.42$ on CHB-MIT, and
replacing factorized attention by a flat encoder costs $0.22$ and $0.53$.

\begin{table}[t]
\centering
\small
\caption{Design variants, single seeds. ``fold'' = band rows
folded into spatial attention; ``strength'' = explicit $|Z|$ feature.}
\label{tab:design}
\resizebox{\textwidth}{!}{%
\begin{tabular}{lcccccccc}
\toprule
Variant & TUSZ & CHB-MIT & TUEV & IIIC & TUAR & ADFD & Siena & Sleep-EDF \\
\midrule
no fold, $d{=}192$ & 0.612 & 0.630 & 0.690 & 0.458 & 0.624 & 0.498 & 0.239 & 0.610 \\
no fold, $d{=}192$, strength & 0.643 & 0.513 & $\mathbf{0.721}$ & 0.448 & 0.626 & 0.467 & $\mathbf{0.398}$ & 0.614 \\
fold, $d{=}192$ (final) & $\mathbf{0.698}$ & $\mathbf{0.716}$ & 0.680 & 0.447 & 0.600 & 0.461 & 0.201 & 0.639 \\
fold, $d{=}256$ & 0.631 & 0.608 & 0.679 & $\mathbf{0.465}$ & $\mathbf{0.657}$ & 0.454 & 0.269 & $\mathbf{0.653}$ \\
\bottomrule
\end{tabular}}
\end{table}

\subsection{The token inside published encoders}
\label{sec:transplant}

Table~\ref{tab:transplant} adds the tokenizer's rows to CBraMod as
described in Section~\ref{sec:encoder}. Against the matched control ---
the same eight extra rows carrying waveform tokens --- the interaction
rows are better on TUEV, TUEP, IIIC and CHB-MIT and worse on TUSZ, so
the effect is attributable to what the rows carry rather than to the
extra channels. Against CBraMod's own tokenizer the transplant is a clear
gain on TUEV ($+0.077$, three standard deviations of the native model) and
on CHB-MIT ($+0.148$, where the native model is unstable across seeds), a
tie on TUEP and IIIC, and a loss on TUSZ.
Table~\ref{tab:hosts} widens the test to two further encoders, with
the tokenizer now \emph{replacing} the host's patch embedding (every
electrode--row pair becomes one host channel, rows are pooled after the
encoder, the head is unchanged). Trained from scratch, LaBraM gains
$0.150$ $\kappa$ on TUEV and REVE $0.241$; on IIIC the gains are
$0.014$ and $0.078$; on CHB-MIT LaBraM gains $0.017$, within its own
seed variance. The three hosts span 4M to 69M parameters and three
encoder designs, and the ordering of the TUEV gains follows how far each
host is from what it can learn from $10^5$ labelled windows on its own.
Adding the same rows to CBraMod's \emph{released} checkpoint is worse
than the checkpoint alone on both corpora tried (appendix): a pretrained
encoder is tuned to its own token distribution, and inserting eight new
channels per electrode disturbs it more than the coupling content
repays. We draw the conclusion the data support --- the token transfers
into encoders trained from scratch, its effect is attributable to the
content of the added rows, and it pays off on epileptiform event
classification --- and not the broad one about released checkpoints.

\begin{table}[t]
\centering
\small
\caption{CroFreMo's rows added to CBraMod's encoder, trained from scratch.
CBraMod's own tokenizer is a three-seed mean; the two transplant arms are
single seeds with identical parameter and token counts.}
\label{tab:transplant}
\begin{tabular}{llccc}
\toprule
Corpus & Metric & CBraMod tokenizer & + waveform rows & + interaction rows \\
\midrule
TUEV & $\kappa$ & $0.564 \pm 0.024$ & $0.596$ & $\mathbf{0.641}$ \\
CHB-MIT & AUC-PR & $0.317 \pm 0.205$ & $0.341$ & $\mathbf{0.465}$ \\
TUEP & AUROC & $0.642 \pm 0.034$ & $0.593$ & $\mathbf{0.650}$ \\
IIIC & $\kappa$ & $0.393 \pm 0.006$ & $0.381$ & $0.396$ \\
TUSZ & AUC-PR & $\mathbf{0.482 \pm 0.052}$ & $0.476$ & $0.384$ \\
\bottomrule
\end{tabular}
\end{table}

\begin{table}[t]
\centering
\small
\caption{CroFreMo's tokenizer in place of the host's own, hosts trained
from scratch. Native rows are three-seed means except REVE from scratch
(single seed); transplants are single seeds. Cohen's $\kappa$; CHB-MIT
AUC-PR.}
\label{tab:hosts}
\resizebox{\textwidth}{!}{%
\begin{tabular}{lcccccc}
\toprule
Host (from scratch) & TUEV native & TUEV + CroFreMo & IIIC native & IIIC + CroFreMo & CHB-MIT native & CHB-MIT + CroFreMo \\
\midrule
CBraMod (4M) & $0.564 \pm 0.024$ & $\mathbf{0.632}$ & $0.393 \pm 0.006$ & $0.396$ & --- & --- \\
LaBraM (5.9M) & $0.372 \pm 0.009$ & $\mathbf{0.522}$ & $0.406 \pm 0.007$ & $0.420$ & $0.360 \pm 0.050$ & $0.377$ \\
REVE (69M) & $0.303$ & $\mathbf{0.544}$ & $0.299$ & $\mathbf{0.377}$ & --- & --- \\
\bottomrule
\end{tabular}}
\end{table}

\subsection{Pretraining}
\label{sec:pretrain-results}

CroFreMo is trained from scratch throughout. A pretraining study on an
earlier design with a dedicated frequency axis --- gains where labels are
scarce or noisy, losses where they are plentiful --- is reported as an
exploratory result in Appendix~\ref{app:pretrain}; we do not build on it.
